\documentclass[11pt]{article}

\usepackage[T1]{fontenc}
\usepackage{amsmath,amssymb}
\usepackage[a4paper,margin=25mm]{geometry}
\usepackage{booktabs}
\usepackage[hidelinks]{hyperref}
\usepackage{microtype}

\newcommand{\ii}{\mathrm{i}}
\newcommand{\kpar}{k_{\parallel}}
\newcommand{\vth}{v_{T\sigma}}
\newcommand{\oc}{\omega_{c\sigma}}
\newcommand{\ld}{\lambda_{D\sigma}}
\newcommand{\Zp}{\mathcal{Z}}
\setlength{\emergencystretch}{3em}

\title{Collisionless Ion Kernels for KIM Forced Periodicity}
\author{}
\date{22 July 2026}

\begin{document}
\maketitle

\begin{abstract}
This note reduces the collisionless magnetized response kernels of Chapter~13
of \texttt{PhD\_thesis.pdf} to the $m_\phi=0$ form implemented by KIM's
forced-periodicity solver.  It records the Fourier convention, finite-Larmor-
radius (FLR) arguments, causal regularization, four implementable charge and
parallel-current kernels, and checks used by the test suite.  The symbolic
reduction and numerical test oracle are reproducible with the accompanying
\texttt{collisionless\_periodic\_kernels.wl} Wolfram Language source.
\end{abstract}

\section{Scope and conventions}

The production periodic model is hybrid: electrons retain the existing
Fokker--Planck response, while every enabled ion species uses the kernels below
when \texttt{ion\_collision\_model=}\allowbreak\texttt{'collisionless'}.
The derivation is
specialized to $m_\phi=0$, consistent with the existing global collisionless
charge implementation.  Species labels are suppressed where unambiguous.

KIM reconstructs radial fields with $\exp(+\ii k_r r)$ and uses the local
two-wavenumber phase
\begin{equation}
  P(k_r,k'_r,r_g)=\exp[-\ii(k_r-k'_r)r_g].
\end{equation}
Each local $G$ below includes the continuous Fourier normalization
$1/(8\pi^2)$ but not $P$.  Thus, the functions supplied to matrix quadrature
are $\widehat G=P G$.  This convention is intentionally the same as the
existing Fokker--Planck Fourier kernels.

The two-wavenumber FLR arguments are
\begin{align}
 b_+ &= \frac{\rho_L^2}{2}\left(2k_s^2+k_r^2+k_r'^2\right),\\
 b_\times &= \rho_L^2\sqrt{k_s^2+k_r^2}\sqrt{k_s^2+k_r'^2},
\end{align}
with overflow-safe scaled Bessel products
\begin{equation}
 S_0=e^{-b_+}I_0(b_\times),\qquad
 S_1=e^{-b_+}I_{-1}(b_\times).
\end{equation}
If the optional global-comparison approximation is enabled, the shared KIM
FLR helper drops $k_s^2$ and takes electrons (but not ions) in the zero-FLR
limit, exactly as for the Fokker--Planck kernels.

\section{Causal prescription}

The collisionless resonance is regularized by a positive user-supplied
$\epsilon$ with dimensions of inverse length:
\begin{equation}
 k_{\mathrm{abs}}=\sqrt{\kpar^2+\epsilon^2},\qquad
 k_{\mathrm{pole}}=\kpar+\ii\epsilon,
\end{equation}
and
\begin{equation}
 z_0=-\frac{\omega_E-\omega}
 {\sqrt{2}\,\vth k_{\mathrm{abs}}},\qquad
 Z_0=\Zp(z_0).
\end{equation}
Here $\Zp$ is the plasma dispersion function.  Keeping $z_0$ real prevents an
unphysical exponentially growing evaluation in its lower half-plane.  Even
charge-response denominators use $k_{\mathrm{abs}}$; signed inverse factors in
$G^{\rho B}$, $G^{j\Phi}$, and $G^{jB}$ use $k_{\mathrm{pole}}$.

\section{\texorpdfstring{$m_\phi=0$}{m-phi=0} reduction}

Setting $m_\phi=0$ in thesis equations (13.53)--(13.55) gives
\begin{align}
 a_0={}&S_0\left[-\frac{\omega_E}{\oc}
 +\frac{k_s\vth^2}{\oc^2}
   \left(A_1+(1+b_+)A_2\right)\right]
 +\frac{k_s\vth^2}{\oc^2}A_2b_\times S_1,\\
 a_1={}&-\frac{k_{\mathrm{abs}}}{\oc}S_0,\\
 a_2={}&\frac{k_s A_2}{2\oc^2}S_0.
\end{align}
The $k_{\mathrm{abs}}$ in $a_1$ is KIM's collisionless continuation of the
signed thesis expression.  It is required for the even $\rho$--$\Phi$
response to recover the same Debye limit on either side of the resonant
surface.

Equations (13.58) and (13.59) reduce to the moments
\begin{align}
 R_0={}&Z_0\left(a_0+\sqrt{2}\vth z_0a_1
                  +2\vth^2z_0^2a_2\right)
       +\sqrt{2}\vth\left(a_1+\sqrt{2}\vth z_0a_2\right),\\
 R_1={}&(z_0Z_0+1)a_0+\sqrt{2}\vth z_0a_1
       +2\vth^2z_0^2a_2+2\vth^2a_2.
\end{align}
The electrostatic kernels from equations (13.63) and (13.67) are therefore
\begin{align}
 G^{\rho\Phi}_\sigma(k_r,k'_r,r_g)
 &=\frac{1}{8\pi^2}
   \frac{\oc}{\sqrt{2}\,\ld^2\vth k_{\mathrm{abs}}}R_0,\label{eq:rhophi}\\
 G^{j\Phi}_\sigma(k_r,k'_r,r_g)
 &=\frac{1}{8\pi^2}
   \frac{\oc}{\ld^2 k_{\mathrm{pole}}}R_1.\label{eq:jphi}
\end{align}

Mathematica reduces the bracket shared by thesis equations (13.151) and
(13.153) to
\begin{equation}
 \mathcal{B}=\frac{A_2}{2}S_0
 +(z_0Z_0+1)\left[
   \left(A_1+A_2(1+b_++z_0^2)\right)S_0+A_2b_\times S_1
 \right].
\end{equation}
The two magnetic-drive kernels are
\begin{align}
 G^{\rho B}_\sigma(k_r,k'_r,r_g)
 &=\frac{1}{8\pi^2}
   \frac{\ii\vth^2}{c\ld^2\oc k_{\mathrm{pole}}}\mathcal{B},\label{eq:rhoB}\\
 G^{jB}_\sigma(k_r,k'_r,r_g)
 &=\frac{1}{8\pi^2}
   \frac{\ii\sqrt{2}\vth^3z_0}
        {c\ld^2\oc k_{\mathrm{pole}}}\mathcal{B}.\label{eq:jB}
\end{align}
No collision frequency occurs in equations~\eqref{eq:rhophi}--\eqref{eq:jB}.
In KIM's Gaussian-cgs convention, the local kernel dimensions implied by these
expressions are summarized in Table~\ref{tab:dimensions}.

\begin{table}[ht]
 \centering
 \caption{Dimensions of the local continuous-Fourier kernels.}
 \label{tab:dimensions}
 \begin{tabular}{@{}cc@{}}
  \toprule
  Kernel & Dimension \\
  \midrule
  $G^{\rho\Phi}$ & $\mathrm{cm}^{-2}$ \\
  $G^{\rho B}$ & $1$ \\
  $G^{j\Phi}$ & $\mathrm{cm}^{-1}\,\mathrm{s}^{-1}$ \\
  $G^{jB}$ & $\mathrm{cm}\,\mathrm{s}^{-1}$ \\
  \bottomrule
 \end{tabular}
\end{table}

\section{Checks and implementation contract}

The symbolic and Fortran tests enforce the following independent checks:
\begin{enumerate}
 \item The four kernels reproduce a nontrivial high-precision Mathematica
       oracle including complex Landau response and off-diagonal FLR coupling.
 \item Equations~\eqref{eq:rhoB} and \eqref{eq:jB} obey
       \begin{equation}
         G^{jB}=\sqrt{2}\vth z_0G^{\rho B}.
       \end{equation}
 \item In a homogeneous static zero-FLR plasma,
       \begin{equation}
         G^{\rho\Phi}=-\frac{1}{8\pi^2\ld^2},\qquad
         G^{\rho B}=G^{j\Phi}=G^{jB}=0.
       \end{equation}
 \item All four kernels remain finite at $\kpar=0$ for $\epsilon>0$ and
       respond to changing $\epsilon$.
 \item Changing the ion Fokker--Planck collision frequency does not change any
       collisionless ion kernel.
 \item The default Fokker--Planck dispatch reproduces the original Fourier
       functions exactly.  Collisionless dispatch equals FP electrons plus
       collisionless ions and honors both species-disable flags.
\end{enumerate}

The current implementation requires positive $\epsilon$, $v_T$, and
$\lambda_D$, nonzero $\omega_c$, and
\texttt{artificial\_debye\_case=}\allowbreak\texttt{0}.  The diagnostic-only artificial Debye
splits are not defined by the compact Chapter~13 reduction and are rejected
instead of being assigned an invented decomposition.  Higher cyclotron
harmonics $m_\phi\ne0$ remain outside the present model.

\end{document}
